% Strategic Management - Per-Slide Research Bullet Notes
% Compile: pdflatex sm-unit-5-6-slides-study.tex

\documentclass[10pt,a4paper]{article}
\usepackage[utf8]{inputenc}
\usepackage[T1]{fontenc}
\usepackage[margin=1.0cm]{geometry}
\usepackage{lmodern}
\usepackage{xcolor}
\usepackage{fancyhdr}
\usepackage{titlesec}
\usepackage{enumitem}

\definecolor{eublue}{RGB}{18,45,104}
\definecolor{notepink}{RGB}{214,36,113}
\definecolor{soft}{RGB}{246,248,252}

\setlength{\parindent}{0pt}
\setlength{\parskip}{0.3em}
\setlist[itemize]{leftmargin=1.2em,itemsep=1.2pt,topsep=1.2pt}

\titleformat{\section}{\normalsize\bfseries\color{eublue}}{}{0pt}{}
\titleformat{\subsection}{\small\bfseries\color{eublue}}{}{0pt}{}

\pagestyle{fancy}
\fancyhf{}
\fancyfoot[L]{\scriptsize Strategic Mgmt | Per-Slide Research Bullets}
\fancyfoot[R]{\scriptsize \thepage}

\newcommand{\slidebullets}[7]{%
\subsection*{#1}
\colorbox{soft}{\parbox{\dimexpr\linewidth-2\fboxsep}{%
\begin{itemize}
  \item \textbf{\textcolor{notepink}{Main point:}} #2
  \item \textbf{\textcolor{notepink}{Key terms:}} #3
  \item \textbf{\textcolor{notepink}{What this means:}} #4
  \item \textbf{\textcolor{notepink}{Why it matters:}} #5
  \item \textbf{\textcolor{notepink}{Class example:}} #6
  \item \textbf{\textcolor{notepink}{Exam angle:}} #7
\end{itemize}
}}\par
}

\begin{document}
\begin{center}
{\Large\bfseries\color{eublue} Strategic Management: Per-Slide Research Bullet Notes}\\
{\small Unit 5 (18 slides) + Unit 6 (13 slides), strict bullet notes per slide}
\end{center}

\section*{Unit 5 - Competitive Strategy (Slides 1-18)}

\slidebullets{Slide 1/18 - Contents}{Unit 5 lays out the logic from competitive advantage basics to life-cycle strategy.}{competitive advantage, cost leadership, differentiation, strategy clock, life-cycle}{The slide is the map for the whole topic and sets the sequence to follow.}{Using this order keeps answers structured and avoids missing key ideas.}{Week 7 class moved in this same order from definitions to application.}{If asked for Unit 5 overview, answer as: concept -> two strategies -> extension -> evolution.}

\slidebullets{Slide 2/18 - Learning Objectives}{The goal is to analyze and apply, not just define, competitive strategy concepts.}{understand, analyze, critically evaluate, suitable strategy}{The professor expects interpretation of when and why strategies work.}{Assessment quality depends on linking theory to context and limits.}{Week 7 discussions compared strategy success across contexts, not memorized labels.}{State objective, mechanism, and boundary condition in every exam answer.}

\slidebullets{Slide 3/18 - Competitive Advantage}{Advantage means better market position and performance through superior value-cost logic.}{competitive advantage, parity, substantial, sustainable, margin}{Advantage is relative to rivals and must last over time to matter.}{This separates real advantage from temporary popularity or noise.}{Starbucks example: brand matters only if it supports durable margins and demand.}{Define advantage, then show source, sustainability, and performance effect.}

\slidebullets{Slide 4/18 - Competitive Strategy}{Competitive strategy is how a firm positions target scope and source of advantage.}{offensive-defensive actions, how to compete, low cost, uniqueness, focus}{Porter matrix forces firms to choose coherent positioning logic.}{Incoherent positions create weak execution and unclear market signal.}{Airline examples show clear cost model vs clear full-service model differences.}{Map the firm on scope and advantage source before judging performance.}

\slidebullets{Slide 5/18 - Cost Leadership Advantage}{Cost leadership exists when comparable offerings are produced at lower cost than rivals.}{cost leadership, industry average, comparable quality, margin, price}{Lower costs allow either stronger margin or lower prices to win share.}{It is a system advantage, not simply ``cheap'' pricing.}{Discount supermarket chains use scale and efficiency to sustain low operating costs.}{Explain both paths: keep price for margin or cut price for volume.}

\slidebullets{Slide 6/18 - Drivers of Cost Leadership}{Cost leadership comes from multiple operational and organizational drivers combined.}{economies of scale, learning, process, design, input costs, utilization}{No single lever is enough; the cost system must reinforce itself.}{This helps identify if low cost is robust or easy to copy.}{Retail logistics plus standardized formats often create cumulative cost edges.}{Name at least three drivers and show how they interact.}

\slidebullets{Slide 7/18 - Adequacy of Cost Leadership}{Cost leadership fits best in price-sensitive and standardized market conditions.}{price sensitivity, commodity, buyer power, scale economies, learning effects}{Context determines strategy fit; low cost is not universally superior.}{Using the wrong context logic causes strategic mismatch and weak returns.}{Supermarket markets with switching and strong buyer power reward low prices.}{Always add: ``This works if customers prioritize price over differentiation.''}

\slidebullets{Slide 8/18 - Risks of Cost Leadership}{Cost advantages erode through innovation shifts, inflation, quality issues, and imitation.}{monitoring costs, innovation shock, inflation, demand shifts, imitation}{Low-cost systems need continuous renewal to avoid decay.}{Risk analysis prevents overconfidence in past efficiency success.}{Ford-style scale logic weakened when demand shifted to new preferences.}{Give one risk and one mitigation action in your answer.}

\slidebullets{Slide 9/18 - Product Differentiation Advantage}{Differentiation creates advantage when customers perceive unique value worth a premium.}{uniqueness perception, price premium, perceived value, differentiation}{Perception is central; features alone do not guarantee advantage.}{Premium pricing works only when value is visible and meaningful to buyers.}{Starbucks ``third place'' value extends beyond coffee product attributes.}{Say what is unique, why customers care, and why they keep paying.}

\slidebullets{Slide 10/18 - Drivers of Differentiation}{Differentiation can come from product, market, firm reputation, and timing factors.}{features, services, intangibles, reputation, legitimacy, time}{Differentiation is multidimensional, including symbolic and relational elements.}{This allows stronger positioning beyond technical product differences.}{Fashion and lifestyle brands often win through identity and experience effects.}{Classify driver type first, then judge imitability and durability.}

\slidebullets{Slide 11/18 - Adequacy of Differentiation}{Differentiation works when valued attributes are salient, rare, and hard to imitate.}{buyer importance, limited crowding, imitation barriers, innovation industries}{Not all uniqueness matters; only relevant and defensible uniqueness pays.}{This is the checklist for judging premium strategy strength.}{Zara-type cases test whether distinctiveness stays meaningful and hard to copy.}{Use a three-part test: relevance, scarcity, and defensibility.}

\slidebullets{Slide 12/18 - Risks of Differentiation}{Differentiation fails when premium exceeds perceived value or rivals copy the offer.}{excess premium, demand contraction, imitation, segment attackers}{Premium strategies are sensitive to value-price balance changes.}{Ignoring affordability and imitation pressures can collapse demand quickly.}{Economic downturns often shift buyers from premium chains to value options.}{Discuss when customers would trade down and what firm should adjust.}

\slidebullets{Slide 13/18 - Critical Learning}{Cost is not price, and strategic labels require evidence not assumptions.}{cost vs price, one cost leader, stuck in the middle, premium logic}{Frameworks are starting points; empirical coherence must be tested.}{This slide prevents formulaic answers and pushes critical thinking.}{Zara discussion showed mid-position can work if system coherence is strong.}{Challenge one label with evidence on capabilities, margins, and customer value.}

\slidebullets{Slide 14/18 - Strategy Clock Intro}{The strategy clock expands Porter by showing more than two viable positions.}{strategy clock, eight positions, price-value combinations, hybrid}{Firms can occupy nuanced positions between pure cost and pure differentiation.}{It improves classification of real-world mixed strategies.}{Both supermarket and airline examples fit better on the clock than binary labels.}{Explain position using both price level and perceived value simultaneously.}

\slidebullets{Slide 15/18 - Strategic Clock Illustration Prompt}{The slide asks applied mapping of eight positions with industry-specific drivers.}{eight positions, air transport, value proposition, drivers, vulnerability}{Good answers require position logic plus operating mechanism per position.}{It trains practical diagnosis, not just framework recitation.}{Airlines and supermarkets provide parallel cases for all eight positions.}{For each position give: offer logic, driver, risk, and likely shift path.}

\slidebullets{Slide 16/18 - Life-Cycle Model Intro}{Competitive strategy must change across industry life-cycle stages.}{emergence, growth, maturity, decline, stage-contingent strategy}{A strategy can be right now and wrong later as conditions evolve.}{Static strategy recommendations often fail because timing is ignored.}{Week 7 discussion emphasized different priorities by stage.}{Always identify stage before recommending competitive moves.}

\slidebullets{Slide 17/18 - Life-Cycle Model Detail}{Each stage has distinct priorities for investment, positioning, and exit choices.}{entry timing, growth investment, shakeout, consolidation, harvest}{Life-cycle logic adds temporal discipline to strategic choices.}{It helps decide when to invest, defend, pivot, or withdraw.}{Consumer tech markets show shakeout signals after rapid growth periods.}{Mention stage signal and action threshold in any recommendation.}

\slidebullets{Slide 18/18 - Wal-Mart China Case Prompt}{The case requires transferability analysis from US advantage to China adaptation.}{sources of success, local adaptation, transferability, forward strategy}{Past success factors must be tested against local constraints and institutions.}{Cross-market fit is the core issue, not simple replication.}{Week 9 Starbucks case used the same logic: source map + renewal risk.}{Answer as: home advantage -> China mismatch -> adapted strategy option.}

\section*{Unit 6 - Corporate Strategy: Development (Slides 1-13)}

\slidebullets{Slide 1/13 - Contents}{Unit 6 moves from business-level competition to corporate scope decisions.}{firm scope, specialization, diversification, vertical integration, restructuring}{The focus changes from ``how to compete'' to ``where to compete.''}{Scope choices shape long-term resource allocation and risk profile.}{Week 8 introduced this as the core corporate-level question.}{Start Unit 6 answers by defining the scope decision first.}

\slidebullets{Slide 2/13 - Learning Objectives}{The unit trains scope-choice reasoning and critical evaluation of development options.}{development, methods, where to compete, restructuring}{Objectives emphasize trade-offs and option appraisal, not list memorization.}{Good answers compare alternatives and justify one path with evidence.}{Week 8 linked each option to conditions, risks, and implementation method.}{Present options, compare fit, then defend chosen strategy.}

\slidebullets{Slide 3/13 - Scope, Development, Growth}{Development (scope change) and growth (size change) are distinct concepts.}{scope modification, quantitative/qualitative change, growth variables, methods}{Corporate direction and execution method must be analyzed separately.}{Confusing ``what'' and ``how'' leads to weak strategic diagnosis.}{Diversification can be pursued via alliances or M\&A with different risk levels.}{Define strategy type first, then choose method and justify.}

\slidebullets{Slide 4/13 - Where to Compete}{Corporate strategy starts with selecting the business domains the firm will enter.}{firm scope, business set, where to compete}{Scope precedes competitive tactics at business level.}{Even good execution fails if the firm chooses poor domains.}{Expansion into adjacent categories must pass scope-fit logic first.}{Answer with scope rationale before operational details.}

\slidebullets{Slide 5/13 - Specialization Divider}{Specialization is concentration on core activity domains.}{specialization, concentration, core business}{This is a depth strategy rather than breadth expansion.}{Concentration can strengthen capability and market position with lower complexity.}{Week 8 framed specialization as low-risk continuity in known domains.}{Differentiate specialization from business-level niche focus in exams.}

\slidebullets{Slide 6/13 - Specialization Reasons}{Specialization is justified by scale, consolidation, incremental sales, and lower risk.}{economies of scale, consolidation, existing products, low risk}{The firm deepens what already works instead of entering new domains.}{This can improve efficiency and defensibility without coordination overload.}{Core-focused firms often grow by improving existing offers and channels.}{When recommending specialization, state why breadth is less attractive now.}

\slidebullets{Slide 7/13 - Diversification Divider}{Diversification is a scope-expansion decision into additional businesses.}{diversification, related, unrelated, scope expansion}{It is a corporate portfolio move, not just growth volume increase.}{Diversification changes complexity, governance, and resource allocation needs.}{Week 8 treated it as a separate path with clear trade-offs.}{State whether diversification is related or unrelated and why.}

\slidebullets{Slide 8/13 - Diversification Reasons and Risks}{Diversification trades potential synergies and risk spread against coordination and cost burdens.}{economies of scope, synergies, risk reduction, inflexibility, entry barriers}{Value creation depends on whether synergies are real and capturable.}{Many diversification failures come from underestimated integration complexity.}{Related diversification usually works better when capabilities transfer clearly.}{Give one synergy source and one execution risk in your answer.}

\slidebullets{Slide 9/13 - Vertical Integration Definition}{Vertical integration extends firm boundaries across multiple value-chain stages.}{make-or-buy boundary, production stages, value chain depth}{This is a control and coordination architecture choice.}{It can reduce dependence but raises organizational demands.}{Manufacturer-supplier integration cases show boundary redesign trade-offs.}{Explain what stage is internalized and why that stage.}

\slidebullets{Slide 10/13 - Vertical Integration Reasons and Risks}{Integration improves control and quality but can raise inflexibility and costs.}{coordination, opportunism, quality guarantee, know-how protection, administrative costs}{Benefits appear when market transactions are risky or strategically sensitive.}{Poorly executed integration destroys flexibility and scale efficiency.}{Using integration to secure technology or quality can justify extra complexity.}{Balance one control gain against one cost/inflexibility downside.}

\slidebullets{Slide 11/13 - Restructuring Definition and Alternatives}{Restructuring redefines scope through refocus, portfolio change, or business re-float.}{redefinition, divestment, core focus, portfolio restructuring, re-float}{It is an active portfolio design tool, not only a crisis reaction.}{Correct option choice determines whether value is recovered or rebuilt.}{Week 8 listed specialization, portfolio restructuring, and business restructuring paths.}{Classify the move first, then explain expected performance effect.}

\slidebullets{Slide 12/13 - Re-float vs Withdrawal}{Underperforming businesses require a choice between recovery measures and exit routes.}{re-float, manager change, debt refinancing, sale, harvest, winding-up}{Decision must match feasibility of turnaround and value-recovery potential.}{Wrong path choice burns resources and delays correction.}{If no recovery path exists, withdrawal alternatives become financially rational.}{State trigger, chosen option, and why alternatives are inferior.}

\slidebullets{Slide 13/13 - Trigger Signals and Causes}{Restructuring begins when businesses miss profitability expectations and show decline signals.}{poor performance, heavy borrowing, inefficient management, demand shift, inertia}{Diagnosis should identify causes before selecting actions.}{Action-first responses fail when root causes are not addressed.}{Week 8 connected trigger signals to cause inventory and option selection.}{Use sequence: signal -> causes -> option -> expected outcome.}

\section*{Cross-Unit Research Bridge}
\begin{itemize}
  \item Unit 5 answers \textbf{how to compete}; Unit 6 answers \textbf{where to compete}; strong case answers combine both levels.
  \item Repeated class inference: execution quality creates success, but adaptation and scope discipline protect durability.
  \item Reusable exam structure: position logic -> capability evidence -> scope choice -> renewal or restructuring implication.
\end{itemize}

\section*{Vocabulary Word Bank}
\begin{itemize}
  \item \textbf{Competitive advantage:} Better relative market position that produces superior performance versus rivals.
  \item \textbf{Competitive parity:} Performance and positioning similar to industry average, without clear advantage.
  \item \textbf{Cost leadership:} Lower cost position than competitors for comparable product/service quality.
  \item \textbf{Differentiation:} Perceived uniqueness that allows a firm to justify a premium.
  \item \textbf{Price premium:} Extra price customers accept due to higher perceived value.
  \item \textbf{Economies of scale:} Lower unit costs as output volume increases.
  \item \textbf{Learning economies:} Cost reduction and quality gains from accumulated experience.
  \item \textbf{Capacity utilization:} Efficient use of fixed assets to spread fixed costs.
  \item \textbf{Buyer power:} Customer ability to pressure prices, terms, or quality demands.
  \item \textbf{Stuck in the middle:} Unclear position with neither strong cost edge nor clear differentiation.
  \item \textbf{Strategy clock:} Eight-position price-value framework extending Porter’s generic strategies.
  \item \textbf{Hybrid strategy:} Combination of reasonable price and meaningful differentiation.
  \item \textbf{Industry life cycle:} Evolution stages (emergence, growth, maturity, decline) shaping strategy fit.
  \item \textbf{Shakeout:} Transition where weaker competitors exit as growth slows.
  \item \textbf{Firm scope:} Set of businesses/activities where a company chooses to compete.
  \item \textbf{Corporate development:} Scope modification through specialization, diversification, integration, or restructuring.
  \item \textbf{Firm growth:} Increase in firm size metrics (sales, assets, outputs, profitability).
  \item \textbf{Specialization:} Concentration on core businesses to deepen position and efficiency.
  \item \textbf{Diversification:} Expansion into additional businesses (related or unrelated).
  \item \textbf{Economies of scope:} Shared resources/capabilities that reduce cost or raise value across businesses.
  \item \textbf{Vertical integration:} Internalization of upstream/downstream value-chain activities.
  \item \textbf{Opportunism risk:} Exposure to self-interested behavior from external partners.
  \item \textbf{Divestment:} Withdrawal from a business through sale, harvest, or closure.
  \item \textbf{Re-float (turnaround):} Actions to recover an underperforming business.
  \item \textbf{Harvest strategy:} Maximize short-term cash flows while reducing long-term investment.
\end{itemize}

\end{document}

